\section{Introduction}

    (Talk about coffee?)

    Tolerance is a known concept in the field of medicine and is defined as the ability of an organism to...
    
    The fundamental principle of tolerance is that it allows the body to adapt to certain stimuli or conditions... 

    When your brain gets used to a certain level of stimulation, it adapts and becomes less responsive to that stimulus. This is why, for example, people who consume caffeine regularly may find that they need to consume more of it to achieve the same stimulating effects.

    In the world of machine learning, it's common to have concepts that are inspired by human physiology, like how neural networks are inspired by the structure of the human brain...

    Of course, In the specific context of Neural Networks, the structures are very simplified and abstracted compared to the complexity of the human brain due to computational constraints and the need for tractability. (It is surprisingly simple since it got invented...?)

    Neurons are just a mathematical function that takes in inputs, applies a simple activation function, and produces an output. They don't have the complex biochemical processes that real neurons have, and they don't have the same kind of plasticity or adaptability that real neurons do.

    It is neither possible nor desirable to create a perfect simulation of the human brain in a neural network, it is also not intelligent to seek complexity for the sake of complexity since we would lose the required performance to train and use the model. However, we can still draw inspiration from the brain and try to incorporate some of its principles into our models in a way that is computationally feasible.

    That's where the concept of tolerance comes in. If we can find a way to incorporate the principle of tolerance into our neural networks, we might be able to create models that are more adaptable and resilient to changes in their environment or input data.

    The idea is simple: information that is new and surprising should be given more weight in the learning process, while information that is familiar and expected should be given less weight. This way, the model can focus on learning from new and informative data, rather than getting stuck in a loop of reinforcing what it already knows.

    ---

    It is known that the number of parameters in a model determines how much information it can retain (Find the article about memorization). Tolerance can help better allocate this limited space and make the model more efficient at learning from new data.

    It can also be understood as a form of negative memory, where the model learns to ignore certain information that it already knows. This can help prevent overfitting and improve generalization to new data. (It is a form of memory since each neuron will save a certain amount of information about the data it has seen)(it is negative since it is the opposite of what we usually think of as memory, which is the ability to retain and recall information, it remembers what it should forget instead of what it should remember)(it can, therefore, improve the models ability in runtime, like when we run the model multiple times)

    \subsection{Background}

    \subsection{Problem}

    \subsection{Objectives}
